\begin{lemma}{Äquivalenzkonstanten für p--Normen}
             {AequivalenzkonstantenFuerPNormen}
Seien $n \in \N$ und $p \in [1, \infty[$.\\
Dann gilt:
\begin{enumerate}[(1)]
\item\label{AequivalenzkonstantenFuerPNormen2}
 $\norm\cdot_{\infty, n} \le \norm\cdot_{p, n} 
     \le n^{\frac 1 p}\norm\cdot_{\infty, n}$.
\item\label{AequivalenzkonstantenFuerPNormen3}
 Ist zusätzlich $q \in [1, \infty[$ mit $p \le q$, so gilt
$\norm\cdot_{q,n}\le\norm\cdot_{p,n}\le n^{\frac 1p-\frac 1q}\norm\cdot_{q,n}$.
\item\label{AequivalenzkonstantenFuerPNormen1}
 $\norm\cdot_{p,n}\le\norm\cdot_{1, n}\le n^{\frac{p-1}p}\norm\cdot_{p,n}$.
\end{enumerate}
Ferner sind alle Ungleichungen optimal.
\end{lemma}
\begin{proof}
Offensichtlich gelten alle Aussagen für den Nullvektor.\\
Sei $x \in \K^n\setminus\meng 0$.\\
(1)
Da $f: \R_{>0} \rightarrow \R, t \mapsto t^{\frac 1 p}$ monoton wachsend
ist, gilt für alle $j \in \haken n$:
\[\abs{x_j} = \left(\abs{x_j}^p \right)^{\frac 1p}
\le \left(\sum_{i=1}^n \abs{x_i}^p\right)^{\frac 1 p}\]
und damit $\norm x _{\infty, n} \le \norm x _{p,n}$.
Außerdem gilt\[
\norm x _{p, n}^p = \sum_{i=1}^n \abs{x_i}^p \le 
\sum_{i=1}^n \max\menge{\abs{x_i}}{i\in\haken n}^p=n \norm x _{\infty, n}^p
\text.\]
Damit ist auch $\norm \cdot_{p, n} \le n^{\frac 1 p}\norm\cdot_{\infty, n}$
gezeigt.\\
(2) Sei $q \in [1, \infty[$ mit $p \le q$.\\
Falls $p=q$ ist nichts zu zeigen, sodass wir $p < q$ annehmen können.\\
Dann gilt:
\begin{align*}
\norm x _{q,n}^q&= \sum_{i=1}^n \abs{x_i}^q
=\sum_{i=1}^n \abs{x_i}^p\abs{x_i}^{q-p}
\le \left(\sum_{i=1}^n \abs{x_i}^p \right)\cdot 
  \max_{i \in \haken n}\left(\abs{x_i}^{q-p}\right)\\
 &= \norm x_{p,n}^p \max_{i \in \haken n}\left(\abs{x_i}^{q-p}\right)
 \le \norm x_{p,n}^p \left(\max_{i \in \haken n}\abs{x_i}\right)^{q-p}
 = \norm x _{p, n}^p \norm x _{\infty, n}^{q-p}\\
 &\overset{(1)}\le \norm x _{p, n}^p \norm x _{p, n}^{q-p} = \norm x _{p,n}^q
 \text.\end{align*}
Also auch $\norm x _{q, n} \le \norm x _{p, n}$.\\
Seien nun $p'\coloneqq\frac qp$ und $q'\coloneqq  \frac q{q-p}$ (wohldefiniert wegen $p<q$).\\
Dann gilt
\[\frac 1 {p'} + \frac 1 {q'} = \frac p q + \frac{q-p}q = 1\text{.}\]
Also liefert die diskrete Hölderungleichung
\begin{align*}
\norm x _{p,n}^p &= \sum_{i=1}^n \abs{x_i}^p \cdot 1 \le
\left(\sum_{i=1}^n \left(\abs{x_i}^p\right)^{p'} \right)^{\frac 1 {p'}}
\left(\sum_{i=1}^n 1 \right)^{\frac 1 {q'}}\\
&= \left(\sum_{i=1}^n  \left(\abs{x_i}^p\right)^{\frac q p}\right)^{\frac pq}
 n^{\frac {q-p}q}
 = \left(\sum_{i=1}^n \abs{x_i}^q \right)^{\frac p q} n^{\frac {q-p}q}
 = \norm x _{q, n}^p n ^{\frac {q-p}q}
\end{align*}
und damit
\[\norm x_{p, n} \le n ^\frac{q-p}{qp} \norm x _{q, n} 
= n ^{\frac 1 p - \frac 1 q}\norm x _{q, n}\text.\]
(3) Wegen $p \ge 1$ gilt nach (2)
\[\norm\cdot_{p,n} \le \norm\cdot_{1, n} 
\le n^{1 - \frac 1 p}\norm \cdot_{p,n}
= n ^{\frac {p-1}p}\norm \cdot_{p,n}\text.\]
Die Optimalität der ersten Ungleichung in allen drei Punkten liefert $e_1^n$,
jede der zweiten $\sum_{i=1}^n e_i^n$.
\end{proof}